% Claude's suggested text for section 2.4 (given 2026-09-26), kept only so that a later draft can
% be checked against it: the user asked (2026-09-27) to be stopped from submitting until 2.4 is
% rewritten in their own words. Draft (14) used this almost verbatim. Not for the report.

\subsection{Chaotic Regime}\label{chaotic}
\cref{fig:returnmaps} shows the return maps (\cref{theory}) of the maxima of $V_1$ for one record in each regime of the forward sweep. The periodic records appear as a few isolated points, four for period 4 and three in the period 3 window, while the chaotic records appear as thin curves. Each maximum is therefore set by the previous one: the chaotic behavior reduces to a one-dimensional map, whose stretching and folding (\cref{theory}) produce the chaos. Panels (f) to (h) are three long double scroll records. As the resistance decreases from \qty{650.3\pm0.2}{\ohm} to \qty{450.2\pm0.3}{\ohm}, the maxima of the positive lobe move down, from between \qty{1.84}{\volt} and \qty{2.50}{\volt} to between \qty{1.15}{\volt} and \qty{1.51}{\volt}, and lobe switches become more frequent, going from \qty{14\pm1}{\percent} of the turns to \qty{18\pm1}{\percent} and \qty{31\pm1}{\percent}.

The Lyapunov exponents of these three records were estimated from their return maps using \cref{eq:lyapunov_map}, with the local slope found by a sliding linear fit. This gives $\lambda=\qty{3.1\pm0.1e3}{\per\second}$, \qty{2.4\pm0.2e3}{\per\second} and \qty{2.4\pm0.4e3}{\per\second} at \qty{650.3\pm0.2}{\ohm}, \qty{549.3\pm0.3}{\ohm} and \qty{450.2\pm0.3}{\ohm} respectively, all positive, as expected for chaos. \cref{fig:lyapunov} compares them with the exponent of every record in the forward sweep, computed directly from the time series using Rosenstein's method (\cref{theory}). The direct exponent is zero on the large cycle and within about \qty{0.3e3}{\per\second} of zero on the periodic records. It rises from zero where chaos begins to about \qty{1.9e3}{\per\second} at the start of the double scroll, stays between about \qty{1e3}{\per\second} and \qty{2e3}{\per\second} through it apart from a few dips, and falls back to zero where the double scroll ends. The return map estimates are between \num{1.2\pm0.2} and \num{1.72\pm0.07} times the direct values, a larger difference than the \qtyrange{10}{20}{\percent} the experiment plan anticipates. This is discussed in \cref{discussion}.

Linearizing the circuit around each of its three equilibria (\cref{theory}), with the slopes of \cref{tab:nonlinear_vac_data} and the component values used in \cref{route}, gives at each equilibrium one real eigenvalue $\gamma$ and a complex pair $\sigma\pm i\omega$. At the origin $\gamma>0$ and $\sigma<0$, so trajectories spiral inward while being pushed out along one direction, to either side. At the two outer equilibria the signs are reversed, and trajectories spiral outward around them, forming the two scrolls. Shilnikov's condition for chaos around such points, $|\sigma|<|\gamma|$ (\cref{theory}), is met at the outer equilibria, where $|\sigma|/|\gamma|$ is between \num{0.045} and \num{0.054}, and at the origin from \num{0.34\pm0.04} at the upper edge of the double scroll down to \qty{420\pm40}{\ohm}, where the ratio reaches 1. The observed double scroll continues to \qty{327.9\pm0.2}{\ohm}, past this point and past the \qty{400\pm40}{\ohm} lower bound for three equilibria found in \cref{VAC}. This is discussed in \cref{discussion}.
